\chapter*{Agradecimientos}

\hspace{12pt}

\begin{flushright}

Toda obra humana es, en el fondo, un pequeño acto de gratitud. Aunque estas páginas lleven mi nombre, sé bien que no me pertenecen por completo. Han sido escritas con el tiempo, el esfuerzo y el afecto de muchas personas que, de una u otra forma, caminaron conmigo. Como las hojas de un árbol que se mecen con el viento sin advertir que forman parte de una misma copa, cada una de ellas aportó algo silencioso y esencial para que este trabajo llegara a existir.

A mis padres, Eduardo y Felicia, primer hogar de mi vida y raíz más profunda de cuanto soy. En ellos conocí la esa forma serena del amor que nada exige y, sin embargo, lo entrega todo. Si aprendí a mirar el mundo con curiosidad y respeto, fue porque ellos me enseñaron que lo más valioso crece despacio y en silencio.

A mi hermano Rony, compañero de los años y de la memoria, cuya nobleza y calidez forman parte inseparable de esa patria íntima llamada familia.

A Jason y Roqui, mis fieles amigos de cuatro patas. En sus pasos tranquilos, en la nobleza de sus ojos y en la sencilla alegría con que acompañaron mis días, encontré una ternura que no necesita palabras. 

A mis tutores, el Dr.~Juan Carlos Cabral, el Dr.~Diego Stalder, el Dr.~José María Gómez, el Dr.~Christian Schaerer y el MSc.~Jesús Núñez, por su guía y generosidad. Deseo expresar un reconocimiento especial al Dr.~Diego Stalder, cuya mentoría serena y constante fue fundamental en el desarrollo de esta tesis. Extiendo también una gratitud particularmente honda al Dr.~Juan Carlos Cabral y al Dr.~Christian Schaerer, con quienes trabajar en el Grupo de Matemática Aplicada y Computación Científica fue una experiencia tan formativa como entrañable.

A mis profesores, el Dr.~Carlos Sauer, el Dr.~Inocencio Ortiz, el Ing.~Juan Zena, el Dr. Fernando Brunetti, al Dr.~Enrique Vargas  y el Ing.~Julio Balbuena, quienes supieron enseñar no solo contenidos, sino una manera honesta y apasionada de aproximarse al saber. 

A Rossana, Ceci, Adrián y Alfredo, con quienes compartí los primeros años de la facultad, cuando todo era todavía descubrimiento y promesa.

A Ángel y Nair, cuya amistad nació en aquellos días extraños en que la distancia separaba los cuerpos, pero no impedía la cercanía del alma. A Nair, en especial, por su fortaleza silenciosa, su calidez y su bondad invariable. Su amistad ha sido una de esas presencias raras y valiosas que iluminan el camino con una luz tranquila.

A Carol, Joseca, Matteo, Lucas, Evelyn, Guille, Meli y Franco, por las tardes de café, las conversaciones compartidas y la compañía que hizo más ligeros los últimos semestres de esta travesía.

A las instituciones que hicieron posible este recorrido: al Programa de Becas del Gobierno del Paraguay, financiado por Itaipú Binacional, por el sostén económico brindado durante mi formación; al  Programa Universitario de Becas para la Investigación, Andrés Borgognon Montero (PUBIABM), por el apoyo otorgado a esta investigación; y al Observatorio Astronómico ``Prof.~Alexis Troche Boggino'', por ofrecer el espacio donde este trabajo pudo encontrar parte de su horizonte.

Finalmente, agradezco a todos aquellos que dejaron en mi vida una palabra oportuna, un gesto amable o una presencia serena. Con los años he comprendido que nada verdaderamente valioso se construye en soledad. Como la luz que atraviesa las hojas al atardecer, toda obra parece simple, pero está hecha de innumerables reflejos.

A todos ellos, mi gratitud más profunda.

\textbf{Carlos María Benítez Cardozo}

\end{flushright}

\bigskip